%% LyX 2.3.4.2 created this file.  For more info, see http://www.lyx.org/.
%% Do not edit unless you really know what you are doing.
\documentclass[xcolor=table]{beamer}
\usepackage{mathpazo}
\renewcommand{\sfdefault}{lmss}
\usepackage[T1]{fontenc}
\usepackage[latin9]{inputenc}
\usepackage{amssymb}
\usepackage{multicol}
\usepackage{siunitx}


\makeatletter
%%%%%%%%%%%%%%%%%%%%%%%%%%%%%% Textclass specific LaTeX commands.
% this default might be overridden by plain title style
\newcommand\makebeamertitle{\frame{\maketitle}}%
% (ERT) argument for the TOC
\AtBeginDocument{%
  \let\origtableofcontents=\tableofcontents
  \def\tableofcontents{\@ifnextchar[{\origtableofcontents}{\gobbletableofcontents}}
  \def\gobbletableofcontents#1{\origtableofcontents}
}

%%%%%%%%%%%%%%%%%%%%%%%%%%%%%% User specified LaTeX commands.
\usetheme{Darmstadt}
%\usetheme{Frankfurt}
% or ...

\setbeamercovered{transparent}
\setbeamersize{text margin left=7mm,text margin right=7mm} 
\makeatother

\usepackage{babel}
\begin{document}
\title[Short Paper Title]{Relativistic Quantum Open Strings}
% \subtitle{Two roads to quantization}

\author[Author]{Harsh Jaluka}
\institute[Universities of Toronto]{Department of Physics\\
University of Toronto}
\date{Final Presentation for PHY472: Intro to String Theory\\ \today}

\makebeamertitle

%\pgfdeclareimage[height=0.5cm]{institution-logo}{institution-logo-filename}

%\logo{\pgfuseimage{institution-logo}}

\AtBeginSection[]{

  \frame<beamer>{ 

    %\frametitle{Outline}   

    \tableofcontents[currentsection] 

  }

}

\addtobeamertemplate{navigation symbols}{}{%
    \usebeamerfont{footline}%
    \usebeamercolor[fg]{footline}%
    \hspace{1em}%
    \insertframenumber/\inserttotalframenumber
}
\setbeamercolor{footline}{fg=black}

%\beamerdefaultoverlayspecification{<+->}
% \begin{frame}{}

% \tableofcontents{}

% \end{frame}

\subsection{Upgrades}
\begin{frame}{Upgrades}
Schrodinger operators: 
\begin{align*}
    X^I(\sigma), ~~x^-_0, ~~\mathcal{P}^{tI}(\sigma), ~~p^+ \tag{12.4}
\end{align*}
Heisenberg operators: 
\begin{align*}
    X^I(\tau, \sigma), ~~x^-_0(\tau), ~~\mathcal{P}^{tI}(\tau, \sigma), ~~p^+(\tau) \tag{12.5} 
\end{align*}
\end{frame}

\subsection{Commutation relations}

\begin{frame}{Commutation relations (CRs)}
Nontrivial:
\begin{align*}
    [X^I(\sigma), \mathcal{P}^{tJ}(\sigma')] &= i \eta^{IJ} \delta(\sigma - \sigma') \tag{12.6}\\
    [x_0^-, p^+] &= -i \tag{12.8} 
\end{align*} 
Trivial:
\begin{align*}
    [X^I(\sigma), X^J(\sigma')] &= [\mathcal{P}^{tI}(\sigma), \mathcal{P}^{tJ}(\sigma')] = 0 \tag{12.7}\\
    [x_0^-, X^I(\sigma)] = [x_0^-, \mathcal{P}^{tI}(\sigma)] &= [p^+, X^I(\sigma)] \\
    &= [p^+, \mathcal{P}^{tI}(\sigma)] = 0 \tag{12.9}
\end{align*}
*Associated CR for Heisenberg operators (a priori) differ only by a $\tau$ dependence in all operators
\end{frame}

\subsection{Hamiltonian}

\begin{frame}{Hamiltonian}
Since $X^+ = 2\alpha' p^+ \tau$, it is easy to guess that*
\begin{align*}
    H &= 2\alpha' p^+ p^- = 2\alpha' p^+ \int_0^\pi d\sigma \mathcal{P}^{\tau-} \tag{12.14}\\
    &= L_0^\perp \tag{9.78, 12.16}
\end{align*}
In all its glory
\begin{align*}
    H(\tau) = \pi \alpha' \int_0^\pi d\sigma \left(\mathcal{P}^{tI} (\tau, \sigma) \mathcal{P}^{tI}(\tau, \sigma) + \frac{X^{I'}(\tau, \sigma) X^{I'}(\tau, \sigma}{(2\pi \alpha')^2}\right) \tag{12.15} 
\end{align*}
\small{*True Hamiltonian is slightly different}
\end{frame}

\begin{frame}{Equations of motion}
Any Heisenberg operator $\xi(\tau, \sigma)$ which arises from a time-independent Schrodinger operator $\xi(\sigma)$ must satisfy 
\begin{align*}
    i \dot{\xi}(\tau, \sigma) = [\xi(\tau, \sigma), H(\tau)] \tag{12.17}
\end{align*}
Hence, 
\begin{align*}
    i \dot{X^I}(\tau, \sigma) &= [X^I(\tau, \sigma), H(\tau)] \\
    &= \vdots \\
    &= 2\pi \alpha' \mathcal{P}^{\tau I}(\tau, \sigma) \tag{12.21}
\end{align*}
coincides with classical equation of motion!
\end{frame}

\begin{frame}{More CRs}
Linear combinations of derivatives $\dot{X}^I \pm {X^I}'$ are simple and useful. 

Nontrivial: 
\begin{align*}
    \left[ (\dot{X}^I \pm {X^I}') (\tau, \sigma), (\dot{X}^J \pm {X^J}') (\tau, \sigma') \right] = \pm 4\pi \alpha' i \eta^{IJ} \frac{d}{d\sigma} \delta(\sigma - \sigma') \tag{12.30}
\end{align*}
Trivial: 
\begin{align*}
    \left[ (\dot{X}^I \pm {X^I}') (\tau, \sigma), (\dot{X}^J \mp {X^J}') (\tau, \sigma') \right] = 0\tag{12.31}
\end{align*}
\end{frame}

\subsection{Recap (and problem)}
\begin{frame}{Recap (and problem)}
    \begin{align*}
    &[X^I(\sigma), \mathcal{P}^{tJ}(\sigma')] = i \eta^{IJ} \delta(\sigma - \sigma') \tag{12.6}\\
    &[x_0^-, p^+] = -i \tag{12.8} \\
    &i \dot{X^I}(\tau, \sigma) = 2\pi \alpha' \mathcal{P}^{\tau I}(\tau, \sigma) \tag{12.21} \\
    &\left[ (\dot{X}^I \pm {X^I}') (\tau, \sigma), (\dot{X}^J \pm {X^J}') (\tau, \sigma') \right] = \pm 4\pi \alpha' i \eta^{IJ} \frac{d}{d\sigma} \delta(\sigma - \sigma') \tag{12.30}
\end{align*}
\begin{itemize}
    \item \textbf{Problem:} We have an infinite set of relations with continuous dependence on $\sigma$. 
    \item Can we get a discrete, enumerable set of relations (creation and annihilation operators...)?
\end{itemize}


\end{frame}

\subsection{Creation and annihilation operators}
\begin{frame}{Creation and annhiliation operators}
Recall 
\begin{align*}
    (\dot{X}^I \pm {X^I}')(\tau, \sigma) = \sqrt{2\alpha' } \sum_{n \in \mathbb{Z}} \alpha_n^I e^{-in(\tau \pm \sigma)} \tag{9.74, 12.33}
\end{align*}
and after a lot of algebra...: \vspace{1em}

Commutation relations between classical modes: 
\begin{align*}
    [\alpha_m^I, \alpha_n^J] = m \eta^{IJ} \delta_{m+n, 0} \tag{12.45}
\end{align*}
Final remaining commutation relation:
% \begin{align*}
%     (\alpha_n^I)^\dagger = \alpha_{-n}^I, ~~~~n \in \mathbb{Z} \tag{12.58}
% \end{align*}
\begin{align*}
    [x_0^I,p^J] = i \eta^{IJ} \tag{12.54}
\end{align*}
\end{frame}

\begin{frame}{Oscillators}
Commutation relations
\begin{align*}
    [a_m^I, a_n^{J \dagger}] = \delta_{m,n} \eta^{IJ} \tag{12.64}
\end{align*}
where 
\begin{align*}
    a_n^I := \frac{\alpha_n^I}{\sqrt{n}}, ~~~~ a_{n}^{I\dagger} := \frac{{\alpha_n^I}^\dagger}{\sqrt{n}} = \frac{\alpha_{-n}^I}{\sqrt{n}} \tag{12.57}
\end{align*}
which gives 
\begin{align*}
    X^I(\tau, \sigma) = x_0^I + 2\alpha' p^I \tau + i \sqrt{2\alpha'} \sum_{n=1}^\infty \left(a_n^I e^{-in\tau} -a_n^{I \dagger} e^{in\tau} \right) \frac{\cos n \sigma}{\sqrt{n}} \tag{12.67} 
\end{align*}
\end{frame}

\subsection{Harmonic oscillators}
\begin{frame}{Strings as harmonic oscillators}
There is an alternative derivation of the 
results we have obtained so far, identifying coefficients of the string Lagrangian with harmonic oscillators with an ever-increasing set of energies. 
\begin{center}
    \textbf{I will not discuss this, in favour of more interesting results later in the chapter}
\end{center}
\end{frame}

\subsection{Transverse Virasoro operators}
\begin{frame}{Transverse Virasoro operators}
We have a mode expansion for $X^I(\tau, \sigma)$ and we know $X^+(\tau, \sigma) = 2\alpha' p^+ \tau$. How about $X^-$? \vspace{1em}

Needed classical transvere Virasoro modes for $X^-$:
\begin{align*}
    L_n^\perp = \frac{1}{2}\sum_{p \in \mathbb{Z}} \alpha_{n-p}^I \alpha_p^I \tag{12.100}
\end{align*}
With $\alpha_n^I$'s now operators...
\begin{center}
    Virasoro \textit{modes} $\to$ Virasoro \textit{operators}
\end{center}
but there is an ordering ambiguity ($[\alpha_{p-n}^I, \alpha_p^I] = (p-n) \delta_{n,0} \implies \alpha_{n-p}^I \alpha_p^I$ don't commute when $n = 0$)
\end{frame}

\subsection{Normal-ordering}
\begin{frame}{Normal ordering to the rescue!}
Normal-ordering: creation on the left, annihilation
on the right 
\begin{align*}
    L_0^\perp &= \frac{1}{2}\alpha_0^I \alpha_0^I + \sum_{p=1}^\infty \alpha_{-p}^I \alpha_p^I = \alpha' p^I p^I + \sum_{p=1}^\infty pa_p^{I\dagger} a_p^I \tag{12.105}
\end{align*}
\textit{without the ordering constant $a$!}
\begin{align*}
    a= \frac{1}{2}(D-2) \sum_{p=1}^\infty p &= -\frac{1}{24} (D-2) \tag{12.107, 12.111}
\end{align*}
where we used the zeta function to evaluate the infinite sum. 


($a$ shifts the mass-squared spectrum and gives rise to massless photon states!)
\end{frame}

\begin{frame}{CRs involving transverse Virasoro operators}
\begin{align*}
    [L_m^\perp, \alpha_n^J] &= -n \alpha_{m+n}^J \tag{12.118} \\
    [L_m^\perp, x_0^I] &= -i \sqrt{2\alpha' }\alpha_m^I \tag{12.119} \\
    [L_m^\perp, L_n^\perp] &= (m-n) L_{m+n}^\perp + \frac{D-2}{12}(m^3 - m) \delta_{m+n, 0} \tag{12.133}
\end{align*}
    
\end{frame}

\begin{frame}{Action on coordinates}
The action of $L_m^\perp$ on coordinates $X^I$ is given by
\begin{align*}
    [L_m^\perp, X^I (\tau, \sigma)] = \xi^\tau_m \dot{X}^I + \xi^\sigma_m {X^I}' \tag{12.136}
\end{align*}
where 
\begin{align*}
    \xi^\tau_m (\tau, \sigma) = -i e^{im\tau} \cos m \sigma, ~~~~\xi^\sigma_m (\tau, \sigma) = e^{im\tau} \sin m \sigma \tag{12.137}
\end{align*}
Indeed, for $m=0$ 
\begin{align*}
    [L_0^\perp, X^I] = -i \partial_\tau X^I 
\end{align*}
the Heisenberg equation of motion for $X^I$.

-- Expected because $L_0^\perp$ is, up to a constant, the string Hamiltonian, and must generate time translations. 
\end{frame}

\subsection{Lorentz generators}
\begin{frame}{What about quantum Lorentz generators?}
Classical Lorentz generators in terms of oscillation modes 
\begin{align*}
    M^{\mu\nu} = x^\mu_0 p^\nu - x_0^\nu p^\mu - i \sum_{n=1}^\infty \frac{1}{n}(\alpha_{-n}^\mu \alpha_n^\nu - \alpha_{-n}^\nu \alpha_{n}^\mu ) \tag{12.147}
\end{align*}
Quantum Lorentz generators are necessary to ensure Lorentz invariance of string theory!
\begin{center}
    In light-cone gauge, $M^{-I}$ turns out to be the most involved quantum Lorentz generator; let's work on this first. 
\end{center}
\end{frame}

\begin{frame}{$M^{-I}$}
Want $M^{-I}$ to be Hermitian and normal-ordered. (Educated) guess 
\begin{align*}
    M^{-I} = x_0^- p^I - \frac{1}{2}(x_0^I p^- + p^- x_0^I) -i \sum_{n=1}^\infty \frac{1}{n}(\alpha_{-n}^- \alpha_n^I - \alpha_{-n}^I \alpha_n^-) \tag{12.150}
\end{align*}
We then find 
\begin{align*}
    [M^{-I}, M^{-J}] = &- \frac{1}{\alpha' (p^+)^2} \sum_{m=1}^\infty (\alpha_{-m}^I \alpha_m^J - \alpha_{-m}^J \alpha_m^I ) \\
    &\cdot \left[ m \left( 1 - \frac{1}{24} (D-2) \right) + \frac{1}{m}\left( \frac{1}{24}(D-2) + a \right) \right] \tag{12.152}
\end{align*}
\end{frame}

\begin{frame}{Dimension of spacetime}
To have $[M^{-I}, M^{-J}] = 0$, we must have 
\begin{align*}
    D = 26, ~~~~ a = -1 \tag{12.155, 12.156}
\end{align*}
``26 dimensions of spacetime!''

The string Hamiltonian now is 
\begin{align*}
    H = L_0^\perp - 1 \tag{12.158}
\end{align*}
This is the true Hamiltonian!
\end{frame}

\subsection{The state space}
\begin{frame}{The state space}
Ground states: 
\begin{align*}
    | p^+, \vec{p}_T \rangle \tag{12.159}
\end{align*}
Annihilated by annihilation operators:
\begin{align*}
    a_n^I | p^+, \vec{p}_T \rangle = 0, ~~~~ n \geq 1, I =2, \cdots, 25 \tag{12.160}
\end{align*}
General basis states created by creation operators: 
\begin{align*}
    | \lambda \rangle = \prod_{n=1}^\infty \prod_{I=2}^{25} \left( a_n^{I\dagger}\right)^{\lambda_{n,I}} |p^+, \vec{p}_T \rangle \tag{12.162}
\end{align*}
leading to an infinite dimensional Hilbert space!
\end{frame}

\begin{frame}{Inner product and time evolution}
The Hilbert space is equipped with an inner product
\begin{align*}
    \langle p'^+, \vec{p'}_T | p^+, \vec{p}_T\rangle = \delta(p'^+ - p^+) \delta(\vec{p'}_T - \vec{p}_T) \tag{12.171}
\end{align*}
Any given state $| \lambda \rangle$ can be time evolved as usual 
\begin{align*}
    \exp\left( -i (L_0^\perp - 1) \tau \right) | \lambda \rangle \tag{12.175}
\end{align*}

\end{frame}

\begin{frame}{Number operator}
Consider the mass-squared operator 
\begin{align*}
    M^2 = \frac{1}{\alpha'} \left( -1 + \sum_{n=1}^\infty n a_n^{I\dagger} a_n^I \right) \tag{12.163}
\end{align*}
We define the numbering operator 
\begin{align*}
    N^\perp = \sum_{n=1}^\infty n a_n^{I\dagger} a_n^I 
\end{align*}
with properties 
\begin{align*}
    [N^\perp, a_n^{I\dagger}] = n a_n^{I\dagger}, ~~~ [N^\perp, a_n^I] = -na_n^I \tag{12.165, 12.166}
\end{align*}
\end{frame}



\begin{frame}{Some particular states}
\begin{center}
    \begin{tabular}{c|c|c|c}
    $N^\perp$ & $|\lambda \rangle$ & $\alpha' M^2$ & Number of states \\ \hline 
    0 & $| p^+, \vec{p}_T \rangle$ & -1 & 1 \\
    1 & $a_1^{I\dagger} |p^+, \vec{p}_T \rangle$ & 0 & 24 \\
    2 & $a_1^{I\dagger} a_1^{J\dagger} |p^+, \vec{p}_T \rangle$, $a_2^{I\dagger} |p^+, \vec{p}_T \rangle$& 1 & 324 
\end{tabular}
\end{center}
Negative masses? Tachyons!
\end{frame}

\begin{frame}{Recap (what have we achieved?)}
Before quantization
    \begin{itemize}
        \item A continuous range of masses with a non-polarized ground state with mass 0. 
        \item No state corresponding to a photon!
    \end{itemize}
After quantization
    \begin{itemize}
        \item Quantized masses with a polarized ground state with mass 0.
        \item Photons!
    \end{itemize}
\end{frame}


\begin{frame}{Equations of motion}
General time-dependent state
\begin{align*}
    | \Psi, \tau \rangle = \int dp^+ d\vec{p_T} \psi_{I_1, \cdots, I_k}(\tau, p^+, \vec{p}_T) a_{n_1}^{I_1 \dagger} \cdots a_{n_k}^{I_k \dagger} | p^+, \vec{p}_T \rangle \tag{12.182}
\end{align*}
Schrodinger equation
\begin{align*}
    &i \frac{\partial }{\partial \tau} |\Psi, \tau \rangle = H | \Psi, \tau \rangle \tag{12.186} \\
    &\implies i \frac{\partial }{\partial \tau}\psi_{I_1, \cdots, I_k} = (\alpha' p^I p^I + N^\perp - 1) \psi_{I_1, \cdots, I_k} \tag{12.188}
\end{align*}
\end{frame}

\subsection{Tachyons}

\begin{frame}{The open string tachyon}
For a free scalar field 
\begin{align*}
    V(\phi) = \frac{1}{2}M^2 \phi^2 
\end{align*}
The equation of motion of a classical, time-dependent, free scalar field 
\begin{align*}
    \partial^2_t \phi - M^2 \phi= 0
\end{align*}
For negative $M^2 = -\beta^2$, as for tachyons, we have 
\begin{align*}
    \phi(t) = A \cosh (\beta t) + B\sinh(\beta t) \tag{12.204} 
\end{align*}
-- Goes to infinity as $t\to \infty$.

-- Instability in open string theory! 
\end{frame}


\subsection{Acknowledgements}
\begin{frame}{Acknowledgements}
\begin{itemize}
    \item Sincere thanks to Professor A.W. Peet for taking their time to go over this text with us!

    \item Thanks to Professor Kushner for taking his time to grade the final presentations!

    \item All equations reference Zwiebach's `A First Course in String Theory'.
\end{itemize}
\end{frame}

\end{document}
